\begin{lemma}\label{lem:supp-simple-equality}
Let $D\in\{9,11,13,15\}$, and let $f(D)$ distinct marked edges form
a strong clique in a simple graph of maximum degree at most $D$.
Unless the marked edges form the five-cycle blow-up with part sizes
$(d,d,d+1,d+1,d)$, where $d=\frac{D-1}{2}$, there is a marked
induced five-cycle with $R\le7$ for $D=9,11$ and $R\le9$ for $D=13,15$.
\end{lemma}

\begin{proof}
Here $\frac{\mu}{D^2}=\frac{5}{4}-\frac{f(D)}{D^2}
=\frac{2D-1}{4D^2}\le\frac{17}{324}$ for $D\ge9$.
If $q_C=0$, the first inequality in \eqref{eq:odd-moment-bounds}
gives $\frac{\mu}{D^2}\ge\frac{\gamma_C}{8}-\epsilon_C
>\frac{17}{324}$. Hence there is an induced red five-cycle.
Choose one maximizing $|B|$. Outside the blow-up construction, the simple-graph equality
check in Appendix~\ref{app:multi-structure} leaves only
$(1,1,2,1,2)$ as the slack vector for $D=9,11$. For $D=13,15$ it
also leaves $(1,1,3,1,3)$ and $(1,2,2,2,2)$, up to the symmetries
of the cycle. The first vector has sum seven and is drawn in
\textbf{Figure~\ref{fig:finite-simple-slack-seven}}. The other two
have sum nine and appear in \textbf{Figure~\ref{fig:finite-nine-slack-nine}}.
\end{proof}

For degrees 11, 13, and 15, we also need to exclude a clique with
exactly one marked copy in excess of its number of marked pairs.

\begin{lemma}\label{lem:single-extra-copy}
For $D\in\{11,13,15\}$, a strong clique with $m=f(D)+1$ marked
copies has $b\ne1$.
\end{lemma}

\begin{proof}
Suppose $b=1$. In the simple support, the $f(D)$ distinct marked
pairs still form a strong clique, since the projection preserves
adjacency between endpoints. It also preserves every marked induced
five-cycle and its five sets $U_i$. In the five-cycle blow-up with class sizes
$(d,d,d+1,d+1,d)$, where $d=\frac{D-1}{2}$, every marked pair has
an endpoint incident with $D$ distinct marked pairs. A second copy
on that pair would give its endpoint degree at least $D+1$ in $G$.

In every other case, \textbf{Lemma~\ref{lem:supp-simple-equality}}
gives a marked induced five-cycle with $R\le7$ at $D=11$, or
$R\le9$ at $D=13,15$. Since the five sets $U_i$ are the same
in the multigraph, these upper bounds
also hold for $R$ there. Each contradicts
\textbf{Lemma~\ref{lem:odd-slack-refinements}}.
\end{proof}
